\chapter{Vecteurs, droites et plans de l'espace}\label{ch:g12:space}

La géométrie en dimension trois devient calculatoire une fois les \omterm{def:g11:vect:vector}{vecteurs}
et les \omterm{def:g10:coordgeom:system}{coordonnées} disponibles~: les droites admettent des représentations
paramétriques, les plans des \omterm{def:g10:algebra:equation}{équations} cartésiennes, et le \omterm{def:g12:space:dot}{produit scalaire}
mesure les angles et les distances. Ce chapitre construit cette boîte à
outils et l'utilise pour résoudre les problèmes classiques d'\omterm{def:g10:numbers:interunion}{intersection}
et de distance.

\section{Vecteurs de l'espace}

Les \omterm{def:g11:vect:vector}{vecteurs} de l'espace obéissent aux mêmes règles que dans le plan~: ils
s'additionnent, se multiplient par un scalaire, et $\vect{AB} = \vect{CD}$
exactement lorsque $ABDC$ est un parallélogramme. Dans un \omterm{def:g10:coordgeom:system}{repère}
$(O; \vec\imath, \vec\jmath, \vec k)$, tout \omterm{def:g11:vect:vector}{vecteur} a des \omterm{def:g10:coordgeom:system}{coordonnées}
$(x, y, z)$.

\begin{definition}[Colinéarité, coplanarité]\label{def:g12:space:collinear}
Deux \omterm{def:g11:vect:vector}{vecteurs} sont \emph{\omterm{def:g11:vect:collinear}{colinéaires}} si l'un est un multiple de l'autre.
Trois \omterm{def:g11:vect:vector}{vecteurs} $\vec u, \vec v, \vec w$ sont \emph{coplanaires}\index{coplanaire}
si l'un d'eux peut s'écrire comme combinaison des deux autres, disons
$\vec w = a\vec u + b\vec v$.
\end{definition}

\begin{proposition}\label{prop:g12:space:basis}
Si $\vec u, \vec v, \vec w$ ne sont pas \omterm{def:g12:space:collinear}{coplanaires}, tout \omterm{def:g11:vect:vector}{vecteur} de
l'espace se décompose de façon unique en $x\vec u + y\vec v + z\vec w$.
\end{proposition}

\begin{proof}[Idée de la preuve]
Par l'extrémité du \omterm{def:g11:vect:vector}{vecteur} à décomposer, mener la droite parallèle à
$\vec w$~; elle rencontre le plan de $\vec u, \vec v$ en un unique point,
qui scinde le \omterm{def:g11:vect:vector}{vecteur} en une composante dans ce plan (uniquement
$x\vec u + y\vec v$, géométrie plane) et une composante selon $\vec w$.
Unicité~: une différence de deux décompositions exprimerait l'un des
\omterm{def:g11:vect:vector}{vecteurs} en \omterm{def:g11:func:function}{fonction} des deux autres, contredisant la non-coplanarité.
\end{proof}

\section{Droites et plans}

\begin{definition}[Représentation paramétrique d'une droite]\label{def:g12:space:line}
La droite passant par $A(x_A, y_A, z_A)$ de \omterm{def:g11:vect:direction}{vecteur directeur}
$\vec u (a, b, c) \neq \vec 0$ est l'ensemble des points
\[
M(x_A + ta,\; y_A + tb,\; z_A + tc), \qquad t \in \R .
\]
\end{definition}

\begin{definition}[Plan]\label{def:g12:space:plane}
Le plan passant par $A$ dirigé par deux \omterm{def:g11:vect:vector}{vecteurs} non \omterm{def:g11:vect:collinear}{colinéaires}
$\vec u, \vec v$ est l'ensemble des points $M$ avec
$\vect{AM} = s\vec u + t\vec v$, $(s, t) \in \R^2$.
\end{definition}

\begin{omfigure}
\begin{tikzpicture}[scale=1.05]
  \draw (-0.5,-0.25) -- (3.4,1.7);
  \draw[-Stealth, omDef, thick] (1,0.5) -- (1.9,0.95)
    node[midway, above left, font=\small] {$\vec u$};
  \fill (0.1,0.05) circle (1.2pt) node[below right, font=\small] {$t=-1$};
  \fill (1,0.5) circle (1.2pt) node[above left, font=\small] {$A$};
  \fill (1.9,0.95) circle (1.2pt) node[below right, font=\small] {$t=1$};
  \fill (2.8,1.4) circle (1.2pt) node[below right, font=\small] {$t=2$};
\end{tikzpicture}
\qquad\quad
\begin{tikzpicture}[scale=1.05]
  \draw (0,0) -- (4,0.5) -- (5.2,1.9) -- (1.2,1.4) -- cycle;
  \node[font=\small] at (4.75,1.6) {$\mathcal P$};
  \coordinate (A) at (1.3,0.55);
  \coordinate (M) at (3.77,1.38);
  \draw[-Stealth, omDef, thick] (A) -- ++(1.31,0.16)
    node[midway, below, font=\small] {$\vec u$};
  \draw[-Stealth, omDef, thick] (A) -- ++(0.42,0.49)
    node[midway, left, font=\small] {$\vec v$};
  \draw[black, dashed, thin] (3.27,0.79) -- (M) -- (1.80,1.14);
  \fill (A) circle (1.2pt) node[below left, font=\small] {$A$};
  \fill (M) circle (1.2pt) node[above, font=\small] {$M$};
\end{tikzpicture}

{\small À gauche~: la droite passant par $A$ de direction $\vec u$, graduée
par le paramètre $t$. À droite~: le plan passant par $A$ dirigé par
$\vec u$ et $\vec v$~; tout point $M$ du plan est atteint comme
$\vect{AM} = s\vec u + t\vec v$.}
\end{omfigure}

\begin{proposition}[Positions relatives]\label{prop:g12:space:positions}
Deux plans distincts sont soit parallèles, soit s'intersectent selon une
droite. Une droite et un plan sont soit parallèles (éventuellement
contenus), soit s'intersectent en un unique point. Deux droites de l'espace
peuvent être \omterm{def:g11:deriv:rate}{sécantes}, strictement parallèles, confondues, ou
\emph{non \omterm{def:g12:space:collinear}{coplanaires}} (gauches).
\end{proposition}

\begin{proof}[Esquisse]
Ce sont des énoncés d'incidence~; chaque analyse de cas se réduit, en
\omterm{def:g10:coordgeom:system}{coordonnées}, à résoudre un \omterm{def:g10:lines:system}{système linéaire} et à compter ses solutions ---
voir \cref{met:g12:space:intersections}.
\end{proof}

\section{Produit scalaire}

\begin{definition}[Produit scalaire]\label{def:g12:space:dot}
Le \emph{produit scalaire}\index{produit scalaire} de $\vec u$ et $\vec v$
est
\[
\vec u \cdot \vec v = \tfrac12\left(\norm{\vec u + \vec v}^2
- \norm{\vec u}^2 - \norm{\vec v}^2\right),
\]
où $\norm{\vec u}$ est la longueur de $\vec u$. Si les deux \omterm{def:g11:vect:vector}{vecteurs} sont
non nuls, $\vec u \cdot \vec v = \norm{\vec u}\,\norm{\vec v}\cos\theta$
où $\theta$ est l'angle entre eux~; et dans un \omterm{def:g10:coordgeom:system}{repère} \emph{orthonormé},
\[
\vec u \cdot \vec v = xx' + yy' + zz' .
\]
Deux \omterm{def:g11:vect:vector}{vecteurs} sont \emph{\omterm{def:g11:scal:orthogonal}{orthogonaux}} si leur produit scalaire vaut $0$.
\end{definition}

\begin{proposition}[Propriétés]\label{prop:g12:space:dotprops}
Le \omterm{def:g12:space:dot}{produit scalaire} est symétrique ($\vec u\cdot\vec v = \vec v\cdot\vec u$),
bilinéaire ($(a\vec u + b\vec v)\cdot \vec w = a\,\vec u\cdot\vec w +
b\,\vec v\cdot \vec w$), et
$\vec u \cdot \vec u = \norm{\vec u}^2 \geq 0$.
\end{proposition}

\begin{proof}
Dans un \omterm{def:g10:coordgeom:system}{repère orthonormé}, les trois propriétés sont immédiates sur la
formule $xx' + yy' + zz'$, qui elle-même suit de la formule définissante et
du calcul pythagoricien des longueurs~:
$\norm{\vec u}^2 = x^2 + y^2 + z^2$.
\end{proof}

\begin{definition}[Vecteur normal]\label{def:g12:space:normal}
Un \emph{vecteur normal}\index{vecteur normal} d'un plan $\mathcal P$ est un
\omterm{def:g11:vect:vector}{vecteur} non nul \omterm{def:g11:scal:orthogonal}{orthogonal} à tout \omterm{def:g11:vect:vector}{vecteur} situé dans $\mathcal P$ (il
suffit qu'il soit \omterm{def:g11:scal:orthogonal}{orthogonal} à deux directions non \omterm{def:g11:vect:collinear}{colinéaires} de
$\mathcal P$).
\end{definition}

\begin{omfigure}
\begin{tikzpicture}[scale=1.1]
  \draw (0,0) -- (4,0.5) -- (5.2,1.9) -- (1.2,1.4) -- cycle;
  \node[font=\small] at (4.75,1.6) {$\mathcal P$};
  \coordinate (H) at (2.6,0.95);
  \coordinate (M0) at (2.86,3.03);
  \draw[-Stealth, omThm, thick] (H) -- (2.73,1.99)
    node[midway, right=1pt, font=\small] {$\vec n$};
  \draw[black, dashed, thin] (2.73,1.99) -- (M0);
  \draw[thin] (2.79,0.974) -- (2.815,1.174) -- (2.625,1.15);
  \draw[black, dashed, thin] (H) -- (1.1,0.55);
  \fill (H) circle (1.2pt) node[below right, font=\small] {$H$};
  \fill (M0) circle (1.2pt) node[above, font=\small] {$M_0$};
  \fill (1.1,0.55) circle (1.2pt) node[below left, font=\small] {$A$};
\end{tikzpicture}

{\small Un \omterm{def:g12:space:normal}{vecteur normal} $\vec n$ du plan $\mathcal P$ est \omterm{def:g11:scal:orthogonal}{orthogonal} à
toute direction de $\mathcal P$. Le point $H$, pied de la perpendiculaire
issue de $M_0$, réalise la distance de $M_0$ au plan.}
\end{omfigure}

\begin{theorem}[Équation cartésienne d'un plan]\label{thm:g12:space:planeeq}
Dans un \omterm{def:g10:coordgeom:system}{repère orthonormé}, le plan passant par $A$ de \omterm{def:g12:space:normal}{vecteur normal}
$\vec n(a, b, c)$ a pour \omterm{def:g10:algebra:equation}{équation}
\[
a x + b y + c z + d = 0,
\]
où $d = -(ax_A + by_A + cz_A)$. Réciproquement, toute telle \omterm{def:g10:algebra:equation}{équation} avec
$(a,b,c) \neq (0,0,0)$ définit un plan de \omterm{def:g12:space:normal}{vecteur normal} $(a, b, c)$.
\end{theorem}

\begin{proof}
$M \in \mathcal P \iff \vect{AM} \perp \vec n \iff
a(x - x_A) + b(y - y_A) + c(z - z_A) = 0$, qui se développe en l'\omterm{def:g10:algebra:equation}{équation}.
Réciproquement, prendre un point solution $A$ quelconque de l'\omterm{def:g10:algebra:equation}{équation}~; le
même calcul à l'envers montre que l'ensemble des solutions est
$\{M : \vect{AM}\cdot\vec n = 0\}$, le plan passant par $A$ de normale
$\vec n$.
\end{proof}

\begin{method}[Intersections en pratique]\label{met:g12:space:intersections}
\begin{itemize}
  \item \emph{Droite $\cap$ plan}~: substituer les \omterm{def:g10:coordgeom:system}{coordonnées} paramétriques
        de la droite dans l'\omterm{def:g10:algebra:equation}{équation} du plan~; résoudre en $t$ (une
        solution~: un point~; $0 = $ non nul~: parallèle~; $0 = 0$~: droite
        contenue).
  \item \emph{Plan $\cap$ plan}~: résoudre le système des deux \omterm{def:g10:algebra:equation}{équations},
        en paramétrant par une coordonnée libre~; la solution est une droite
        (ou les plans sont parallèles lorsque les \omterm{def:g12:space:normal}{vecteurs normaux} sont
        \omterm{def:g11:vect:collinear}{colinéaires}).
  \item \emph{Contrôles d'orthogonalité}~: droite $\perp$ plan ssi sa
        direction est \omterm{def:g11:vect:collinear}{colinéaire} à la normale~; deux plans sont
        perpendiculaires ssi leurs normales sont \omterm{def:g11:scal:orthogonal}{orthogonales}.
\end{itemize}
\end{method}

\begin{proposition}[Distance d'un point à un plan]\label{prop:g12:space:distance}
\index{distance!point à un plan}
Dans un \omterm{def:g10:coordgeom:system}{repère orthonormé}, la distance de $M_0(x_0, y_0, z_0)$ au plan
$\mathcal P : ax + by + cz + d = 0$ est
\[
\operatorname{dist}(M_0, \mathcal P)
= \frac{\abs{ax_0 + by_0 + cz_0 + d}}{\sqrt{a^2 + b^2 + c^2}}.
\]
\end{proposition}

\begin{proof}
Soit $H$ la projection \omterm{def:g11:scal:orthogonal}{orthogonale} de $M_0$ sur $\mathcal P$~: le pied de
la perpendiculaire, donc $\vect{HM_0}$ est \omterm{def:g11:vect:collinear}{colinéaire} à la normale unitaire
$\frac{\vec n}{\norm{\vec n}}$, et la distance est
$\abs{\vect{HM_0}\cdot \frac{\vec n}{\norm{\vec n}}}$. Pour tout
$H \in \mathcal P$,
$\vect{HM_0}\cdot\vec n = ax_0 + by_0 + cz_0 - (ax_H + by_H + cz_H)
= ax_0 + by_0 + cz_0 + d$ (en utilisant l'\omterm{def:g10:algebra:equation}{équation} en $H$), d'où la formule
après division par $\norm{\vec n} = \sqrt{a^2+b^2+c^2}$.
\end{proof}

\begin{example}\label{ex:g12:space:distance}
Distance de l'origine au plan $x + 2y + 2z - 6 = 0$~:
$\frac{\abs{-6}}{\sqrt{1 + 4 + 4}} = \frac63 = 2$.
\end{example}

\section{Exercices}

Dans tous les exercices le \omterm{def:g10:coordgeom:system}{repère} est orthonormé.

\begin{exercise}[$\star$]\label{exo:g12:space:1}
Soient $A(1, 0, 2)$, $B(3, 1, 1)$ et $C(2, -1, 3)$. Les points $A$, $B$,
$C$ sont-ils alignés~? Donner une représentation paramétrique de la droite
$(AB)$.
\end{exercise}

\begin{exercise}[$\star$]\label{exo:g12:space:2}
Déterminer une \omterm{def:g10:algebra:equation}{équation} cartésienne du plan passant par $A(1, 1, 0)$ de
\omterm{def:g12:space:normal}{vecteur normal} $\vec n(2, -1, 3)$. Le point $B(0, 2, 1)$ lui
appartient-il~?
\end{exercise}

\begin{exercise}[$\star$]\label{exo:g12:space:3}
Calculer l'angle en $A$ (au degré le plus proche) du triangle de sommets
$A(0,0,0)$, $B(1, 1, 0)$ et $C(1, 0, 1)$.
\end{exercise}

\begin{exercise}[$\star\star$]\label{exo:g12:space:4}
On considère la droite $\mathcal D$ passant par $A(1, 2, 0)$ de direction
$\vec u(1, -1, 2)$, et le plan $\mathcal P : 2x + y - z + 1 = 0$.
Déterminer $\mathcal D \cap \mathcal P$.
\end{exercise}

\begin{exercise}[$\star\star$]\label{exo:g12:space:5}
Soient $\mathcal P : x - y + 2z = 1$ et $\mathcal Q : 2x + y + z = 4$.
Montrer que $\mathcal P$ et $\mathcal Q$ s'intersectent selon une droite et
en donner une représentation paramétrique.
\end{exercise}

\begin{exercise}[$\star\star$]\label{exo:g12:space:6}
Montrer que les droites
\[
\mathcal D_1 : (x, y, z) = (1 + t,\; 2t,\; 3 - t), \qquad
\mathcal D_2 : (x, y, z) = (2 + s,\; 1 + s,\; 1 + 2s)
\]
sont non \omterm{def:g12:space:collinear}{coplanaires} (gauches).
\end{exercise}

\begin{exercise}[$\star\star$]\label{exo:g12:space:7}
Soit $\mathcal P : 2x - y + 2z - 3 = 0$.
\begin{enumerate}
  \item Calculer la distance de $M_0(3, 1, 1)$ à $\mathcal P$.
  \item Déterminer les \omterm{def:g10:coordgeom:system}{coordonnées} de la projection \omterm{def:g11:scal:orthogonal}{orthogonale} $H$ de
        $M_0$ sur $\mathcal P$, et vérifier la distance $M_0H$.
\end{enumerate}
\end{exercise}

\begin{exercise}[$\star\star\star$]\label{exo:g12:space:8}
Le cube $ABCDEFGH$ a pour arête $1$~; placer les \omterm{def:g10:coordgeom:system}{coordonnées} de sorte que
$A(0,0,0)$, $B(1,0,0)$, $D(0,1,0)$, $E(0,0,1)$ (avec
$C = B + D$, $F = B + E$, $G = B + D + E$, $H = D + E$ comme \omterm{def:g11:vect:vector}{vecteurs} issus
de $A$).
\begin{enumerate}
  \item Montrer que la diagonale $(AG)$ est \omterm{def:g11:scal:orthogonal}{orthogonale} au plan $(BDE)$.
  \item Calculer la distance de $A$ au plan $(BDE)$, et le point où $(AG)$
        le perce.
\end{enumerate}
\end{exercise}

\begin{exercise}[$\star\star\star$]\label{exo:g12:space:9}
\emph{(Volume d'un tétraèdre.)} Soient $A(1,1,1)$, $B(2,1,0)$, $C(0,2,1)$
et $D(2,2,2)$.
\begin{enumerate}
  \item Vérifier que $\vect{AB}$, $\vect{AC}$, $\vect{AD}$ ne sont pas
        \omterm{def:g12:space:collinear}{coplanaires} (les points forment un vrai tétraèdre).
  \item Trouver un \omterm{def:g11:vect:vector}{vecteur} $\vec n$ \omterm{def:g11:scal:orthogonal}{orthogonal} à la fois à $\vect{BC}$ et
        $\vect{BD}$, et en déduire une \omterm{def:g10:algebra:equation}{équation} cartésienne du plan
        $(BCD)$.
  \item Calculer la distance de $A$ au plan $(BCD)$ et l'aire du triangle
        $BCD$, en utilisant la formule
        $\text{Aire} = \frac12\sqrt{\norm{\vec u}^2\norm{\vec v}^2 -
        (\vec u \cdot \vec v)^2}$ pour le triangle engendré par
        $\vec u, \vec v$.
  \item En déduire le volume du tétraèdre
        ($V = \frac13 \times \text{aire de base} \times \text{hauteur}$).
\end{enumerate}
\end{exercise}
